\section{Empirical Status and Direct Validation}
\label{sec:experiments}

The current evidence has two roles: it shows that finite-generation
instability is visible in the existing pipeline, and it identifies the
experiment still needed to test the calibrated buffer.  We keep those roles
separate.  In particular, none of the results below is presented as a completed
evaluation of Theorem~\ref{thm:buffered-main}.

\subsection{What the tracked artifacts establish}
\label{sec:experiments-existing}

The repository contains tracked generation files for AIME 2024, AIME 2025,
GPQA-Diamond, and MATH-500 \citep{rein2023gpqaag,hendrycks2021measuringmp}.
The reproducible analysis in
\texttt{zhi\_code/src/pacbayes\_boinf.py} compares the original weight-fitting
MILP with a bootstrap-stability plus max-margin variant.  Averaged over the four
benchmarks, the absolute train--test gap changes from \(0.0831\) to \(0.0717\),
or \(13.67\%\) relative closure.  Across five stored cross-benchmark
directions, the mean closure is \(26.00\%\), although the individual directions
are heterogeneous and one is negative.  The per-benchmark values, figures, and
artifact paths are reported in Appendix~\ref{app:legacy-real}; every stored
in-domain, transfer, model-count, and PAC--Bayes value is transcribed in
Appendix~\ref{app:legacy-full-data}.

These diagnostics establish that the code path is executable and that
regularizing a weight fit can change finite-sample generalization.  They do not
isolate the theorem's mechanism: the implemented variant combines bootstrap
stability with a numerical margin, whereas buffered ERM uses the explicit
uncertainty radius \(\varepsilon_N(\delta)\).  The retained PAC--Bayes values
are likewise diagnostics, not certificates for the deterministic buffered
predictor.

The earlier draft also reported a small synthetic mechanism check.  Its four
rows show the mass of the \(N^{-1/2}\) boundary band decreasing with \(N\),
while sign flips concentrate inside that band; a second column tracks the
selection bias from optimizing noisy empirical accuracy.  Because the original
per-replicate files are unavailable, Appendix~\ref{app:legacy-synthetic}
labels these values as legacy reports rather than a fresh rerun.

\subsection{The direct buffered-ERM experiment}
\label{sec:experiments-required}

A direct test should hold the model set and answer extractor fixed, repeatedly
split benchmark questions into training and test sets, and independently
subsample both axes: \(Q\) labeled training questions and \(N\) generations per
model and question.  On every split, it should compare uniform voting, raw ERM
with zero buffer, buffered ERM with
\(\varepsilon_N(\delta)=\sqrt{2\log(8KAQ/\delta)/N}\), and the existing
stability-plus-margin baseline.  The multiplicative constant in the buffer may
be selected on a validation split, but the test split must not influence it.

The primary outcomes are clean held-out error under high-generation
frequencies, end-to-end finite-generation held-out error, excess error relative
to a high-generation oracle fit, and the train--test gap.  Reporting the two
held-out errors separately prevents prediction-time sampling noise from being
mistaken for weight-estimation error.  Mechanism checks should also report the
empirical mass of the positive comparator band
\(0<M_q(w^\star)\le2\varepsilon_N\), margin-sign flips inside and outside that
band, and sensitivity to \(Q\), \(N\), \(K\), and the confidence level.  Repeated
splits or a paired bootstrap over questions should provide uncertainty
intervals.  This design distinguishes the \(Q\)-generalization term from the
\(N\)-resolution term instead of inferring both from one aggregate accuracy
curve.

\subsection{Implementation boundary}
\label{sec:experiments-implementation}

The existing solver in \texttt{src/tte/milp.py} already exposes the relevant
constraint through \texttt{strict\_margin}.  Passing
\(\varepsilon_N(\delta)\) implements Eq.~(\ref{eq:buffered-milp}); scripts that
pass \(10^{-4}\) use only a numerical tie-breaking tolerance.  Promoting the
planned experiment to a completed result therefore requires saved split
indices, sampled-generation indices, fitted weights, per-split metrics, and
the aggregate table or figure.  Until those artifacts exist, the theoretical
guarantee and the legacy diagnostics remain intentionally distinct.

The current matrix builder also forms each answer set from the union of answers
observed in the same finite generations.  A theorem-aligned run must instead
fix the answer normalization and candidate support before fitting, or define a
residual category whose treatment preserves the strongest-contender rule.
Otherwise an unobserved wrong answer can disappear from the empirical
\(\arg\max\), and the fixed-\(A\) concentration argument does not directly
describe the implemented decision rule.
